\documentclass[11pt,a4paper]{article}
\usepackage[margin=22mm]{geometry}
\usepackage{amsmath,amssymb}
\usepackage{fontspec}
\usepackage{polyglossia}
\setdefaultlanguage{russian}
\setmainfont{DejaVu Serif}
\usepackage[unicode,hidelinks]{hyperref}

\allowdisplaybreaks[1]
\raggedright
\setlength{\parskip}{0.35em}
\newcommand{\dd}{\mathop{}\!\mathrm{d}}
\newcommand{\avg}[1]{\langle #1\rangle}
\newcommand{\ket}[1]{\lvert #1\rangle}
\newcommand{\bra}[1]{\langle #1\rvert}
\newcommand{\braket}[2]{\langle #1\mid #2\rangle}
\DeclareMathOperator{\Var}{Var}

\newcommand{\solutiontitle}[4]{%
  \hypersetup{pdftitle={Квантовая механика 1. Решения ДЗ №#1}}%
  \begin{center}
    {\small Квантовая механика 1 \enspace\textperiodcentered\enspace Семестр 1\par}
    \vspace{1ex}
    {\LARGE\bfseries Решения домашнего задания №#1\par}
    \vspace{0.5ex}
    {\large #2\par}
  \end{center}
  \noindent{\small Условия приведены в файле «#3», #4.}\par\medskip
}

\begin{document}
\solutiontitle{7}{Бесконечная яма и поток вероятности}{ДЗ2026(1)\_07 (1).pdf}{с.~1}

\section{Импульсное представление стационарных состояний}

Внутри ямы решение уравнения \(-\hbar^2 u''/(2m)=Eu\)
с \(u(0)=u(a)=0\) пропорционально \(\sin(k_nx)\),
где \(k_n=n\pi/a,\ n\ge1\). Норма квадрата синуса на
\([0,a]\) равна \(a/2\), так что \(u_n(x)=\sqrt{2/a}\sin(k_nx)\)
внутри и ноль снаружи. Энергия \(E_n=\hbar^2k_n^2/(2m)\).
При \(q=p/\hbar\) представление синуса через экспоненты даёт
\begin{align*}
 a_n(p)&=\frac1{\sqrt{\pi a\hbar}}\int_0^a
                         \sin(k_nx)e^{-iqx}\,\dd x\\
       &=\frac1{\sqrt{\pi a\hbar}}\,
           \frac{k_n[1-(-1)^n e^{-iqa}]}{k_n^2-q^2}.
\end{align*}
При \(q=\pm k_n\) интеграл равен \(\mp ia/2\) до умножения
на \(1/\sqrt{\pi a\hbar}\). Это устранимые особенности.
Нормировка по \(\dd p\) следует из Планшереля.
Здесь измерение импульса описывается преобразованием Фурье
продолженной нулём функции на всей оси.

\section{Суперпозиция первых двух уровней}

Обозначим \(E_1=\pi^2\hbar^2/(2ma^2)\) и
\(\Omega=(E_2-E_1)/\hbar=3E_1/\hbar\).

\textbf{(a, b)} Разложение в нормированном базисе имеет
коэффициенты \(1/\sqrt5\) и \(2i/\sqrt5\) при \(t=0\).
Стационарные фазы и вероятности уровней дают
\[
 \Psi(x,t)=\frac{u_1(x)e^{-iE_1t/\hbar}
                    +2i\,u_2(x)e^{-i4E_1t/\hbar}}{\sqrt5},
 \qquad \avg{E}=\frac{E_1+4E_2}{5}=\frac{17}{5}E_1 .
\]

\textbf{(c)} Пусть \(z=c_1^*c_2=(2i/5)e^{-i\Omega t}\).
Для каждого эрмитова \(B\)
\(\avg{B}=(B_{11}+4B_{22})/5+2\operatorname{Re}(zB_{12})\).
Тождество \(2\sin u\sin v=\cos(u-v)-\cos(u+v)\) сводит интегралы
для \(x\) и \(x^2\) к интегралам по \(\cos(l\pi x/a)\).
При \(k=l\pi/a\ne0\) интегрирование по частям даёт
\(\int_0^a x\cos(kx)\dd x=[(-1)^l-1]/k^2\) и
\(\int_0^a x^2\cos(kx)\dd x=2a(-1)^l/k^2\).
Постоянные части в \(u_n^2\) и разность косинусов в
\(u_1u_2\) дают
\begin{align*}
 x_{nn}&=a/2,& (x^2)_{nn}&=a^2\left(\frac13-\frac{1}{2\pi^2 n^2}\right),\\
 x_{12}&=-\frac{16a}{9\pi^2},& (x^2)_{12}&=-\frac{16a^2}{9\pi^2}.
\end{align*}
Введём \(d=64a/(45\pi^2)\). Подстановка даёт
\begin{align*}
 \avg{x}&=a/2-d\sin\Omega t,\\
 \avg{x^2}&=a^2\left(\frac13-\frac1{5\pi^2}\right)-ad\sin\Omega t,\\
 \Var(x)&=a^2\left(\frac1{12}-\frac1{5\pi^2}\right)-d^2\sin^2\Omega t .
\end{align*}
Для \(p\) диагональные элементы нулевые, поскольку
\(\int_0^a u_n u_n'\dd x=[u_n^2/2]_0^a=0\). Из
\[
 p_{12}=-i\hbar\frac{4\pi}{a^2}
             \int_0^a\sin(\pi x/a)\cos(2\pi x/a)\dd x
        =\frac{8i\hbar}{3a}
\]
получаем \(\avg{p}=-32\hbar\cos\Omega t/(15a)\).
Второй момент в преобразовании Фурье равен
\(\hbar^2\int_0^a|\partial_x\Psi|^2\dd x\).
Граничные значения равны нулю, поэтому при интегрировании по
частям \(\int_0^a u_i'u_j'\dd x=k_j^2\delta_{ij}\).
В итоге
\[
 \avg{p^2}=\frac{17}{5}\left(\frac{\pi\hbar}{a}\right)^2,\qquad
 \Var(p)=\frac{17}{5}\left(\frac{\pi\hbar}{a}\right)^2
                    -\left(\frac{32\hbar}{15a}\right)^2\cos^2\Omega t .
\]
Проверка первого момента даёт \(m\,\dd\avg{x}/\dd t=\avg{p}\).

\textbf{(d)} Ток равен
\((\hbar/m)\operatorname{Im}(z)(u_1u_2'-u_2u_1')\).
При \(u=\pi x/a\) разность равна
\[
 \frac{4\pi}{a^2}\sin u\cos 2u-\frac{2\pi}{a^2}\sin 2u\cos u
 =-\frac{4\pi}{a^2}\sin^3u.
\]
Поэтому \(j(x,t)=-8\pi\hbar\,\sin^3(\pi x/a)\cos\Omega t/(5ma^2)\)
в яме, снаружи ток равен нулю.

\textbf{(e)} Каждое \(u_n^2\) по отдельности имеет интеграл \(1/2\)
на левой половине. Перекрёстный интеграл равен
\[
 \int_0^{a/2}u_1u_2\,\dd x
 =\frac1\pi[\sin u-\tfrac13\sin 3u]_0^{\pi/2}=\frac4{3\pi}.
\]
Следовательно \(\Pr(x<a/2)=1/2+(16/(15\pi))\sin\Omega t\).
Его производная равна \(-j(a/2,t)\), как следует из уравнения
непрерывности.

\section{Симметричная яма и квадратичное начальное состояние}

В переменной \(y=x+a/2\) границы равны \(0,a\), поэтому
\[
 v_n(x)=\sqrt{2/a}\sin\!\frac{n\pi(x+a/2)}a,\qquad
 E_n=\frac{\pi^2\hbar^2 n^2}{2ma^2}\quad (n=1,2,\ldots).
\]
Стационарное состояние в момент \(t\) есть \(v_n e^{-iE_nt/\hbar}\),
снаружи оно нулевое. Синусы полны по теореме о
рядах Фурье для нечётного продолжения функции с \([0,a]\).
Квадратичный профиль равен \(A(y^2-ay)\).
Из \(\int_0^a(y^2-ay)^2\dd y=a^5/30\) берём \(A=\sqrt{30}/a^{5/2}\).
Коэффициенты получаются двумя интегрированиями по частям, при
которых профиль и синус обнуляются на концах. Получаем
\begin{align*}
 \int_0^a(y^2-ay)\sin(k_n y)\,\dd y
   &=-\frac{2}{k_n^2}\int_0^a\sin(k_n y)\,\dd y
     =-\frac{2[1-(-1)^n]}{k_n^3},\\
 c_n&=-\frac{2\sqrt{60}}{\pi^3 n^3}\,[1-(-1)^n].
\end{align*}
Доступны уровни с нечётным \(n\), для которых
\(\Pr(E_n)=960/(\pi^6n^6)\). Для чётных \(n\) вероятность нулевая.
Волновая функция \(\Psi(x,t)=\sum_{n\ge1}c_n v_n(x)e^{-iE_nt/\hbar}\).
Квадратичная форма гамильтониана в начальный момент даёт
\[
 \avg{E}=\frac{\hbar^2}{2m}\int_{-a/2}^{a/2}|2Ax|^2\,\dd x
        =\frac{\hbar^2}{2m}\frac{4A^2a^3}{12}
        =\frac{5\hbar^2}{ma^2}.
\]
Энергия сохраняется, так как фазы коэффициентов не меняют
вероятностей измерения.

\section{Токи для заданных волновых функций}

Используем \(\mathbf j=(\hbar/m)\operatorname{Im}(\Psi^*\nabla\Psi)\),
эквивалентную формуле из условия.

\textbf{(a)} У \(\cos kx\) и его производной вещественные значения,
а множитель \(Ce^{-iEt/\hbar}\) имеет одинаковую пространственную
фазу всюду. Поэтому \(j=0\).

\textbf{(b)} Положим \(z=C_1^*C_2e^{-i(3kx+3Et/\hbar)}\).
Производные первой и второй волн дают \(ik\) и \(-2ik\).
Самодействия дают \(ik|C_1|^2-2ik|C_2|^2\), перекрёстные члены
дают \(-2ikz+ikz^*\). Итак,
\[
 j=\frac{\hbar k}{m}\bigl(|C_1|^2-2|C_2|^2-\operatorname{Re}z\bigr).
\]

\textbf{(c)} Для плоской волны
\(\nabla\Psi=(i\mathbf p/\hbar)\Psi\), значит
\(\mathbf j=|C|^2\mathbf p/m\).

\textbf{(d)} Для сферической волны с амплитудой \(C/r\)
\(\nabla\Psi=(-1/r+ip/\hbar)\Psi\,\mathbf e_r\) при \(r>0\).
Следовательно \(\mathbf j=|C|^2p\,\mathbf e_r/(mr^2)\).
Поток через сферу радиуса \(r>0\) равен \(4\pi |C|^2p/m\).

\textbf{(e)} Для вещественных \(A,S\) имеем
\(\Psi^*\nabla\Psi=A\nabla A+iA^2\nabla S\).
Отсюда \(\mathbf j=(\hbar/m) A^2\nabla S\).

\end{document}
